% GENERATED FILE. Edit canonical lessons, then run npm run generate:books.
% canonical-source-hash: fnv1a64:b4db8a44d4e735a2
% canonical-lessons: MR-C470-labad, MR-C470-dhongi, MR-C470-bhola, MR-C470-niragas, MR-C470-thakavne

\chapter{Trust and Trickery}
\label{ch:mr-trust-and-trickery}
\hlchaptermodality{470}

\begin{chapteropening}
\textbf{By the end of this chapter:} \emph{I can describe a sly, hypocritical, naive or innocent person.}

Complete the last lesson of chapter 470: I can describe a sly, hypocritical, naive or innocent person.
\end{chapteropening}

\section[\texorpdfstring{\mr{लबाड} (labāḍ)}{लबाड (labāḍ)}]{\mr{लबाड} (labāḍ) — sly; dishonest}

\label{lesson:MR-C470-labad}



\begin{quote}
Before the new one: say the Marathi for fierce, then the Marathi for deceitful.
\end{quote}

\subsection*{You'll want to know: \mr{लबाड}}
\textbf{\mr{लबाड}} — \emph{labāḍ} — \textquotedblleft{}sly; dishonest\textquotedblright{}.

\begin{quote}
\mr{तो} \mr{लबाड} \mr{दुकानदार} \mr{कमी} \mr{वजन} \mr{देतो}.
\end{quote}

\begin{quote}
\emph{to labāḍ dukāndār kamī vajan deto.}
\end{quote}

\begin{quote}
That sly shopkeeper gives short weight.
\end{quote}

Unchanging; also a noun, a cheat. \mr{लबाडी}, trickery.

\begin{cousinweb}
\textbf{\mr{लबाड}} (\emph{labāḍ}) is inherited: CDIAL lists Marathi \emph{labāḍ}, \textquotedblleft{}false, untrue\textquotedblright{}, under the reconstructed \emph{*lappha}, \textquotedblleft{}defective\textquotedblright{} (CDIAL 10943).
\end{cousinweb}

\begin{grammarlens}[title={What you've built}]
One more word for describing what you lived and felt.
\end{grammarlens}

\subsection*{Your turn}
\begin{itemize}
\raggedright
  \item \emph{Say it:} \emph{labāḍ}
  \item \emph{Say it:} \emph{labāḍ}, once more
  \item \emph{Say it:} \emph{to labāḍ dukāndār kamī vajan deto.} — That sly shopkeeper gives short weight.
  \item \emph{Recall:} say the Marathi for fierce, then the Marathi for deceitful, then say \emph{labāḍ} again
\end{itemize}

\begin{tcolorbox}[breakable,colback=teal!4,colframe=teal!35!black,title={Before you move on}]
What does \mr{लबाड} mean? (\textquotedblleft{}Sly; dishonest\textquotedblright{}.) Say it once more.
\end{tcolorbox}
\section[\texorpdfstring{\mr{ढोंगी} (ḍhoṅgī)}{ढोंगी (ḍhoṅgī)}]{\mr{ढोंगी} (ḍhoṅgī) — hypocritical}

\label{lesson:MR-C470-dhongi}



\begin{quote}
Before the new one: say the Marathi for deceitful, then the Marathi for sly.
\end{quote}

\subsection*{You'll want to know: \mr{ढोंगी}}
\textbf{\mr{ढोंगी}} — \emph{ḍhoṅgī} — \textquotedblleft{}hypocritical\textquotedblright{}.

\begin{quote}
\mr{बाहेर} \mr{गोड} \mr{बोलणारा} \mr{तो} \mr{माणूस} \mr{ढोंगी} \mr{आहे}.
\end{quote}

\begin{quote}
\emph{bāher goḍ bolṇārā to māṇūs ḍhoṅgī āhe.}
\end{quote}

\begin{quote}
That man who talks sweetly in public is a hypocrite.
\end{quote}

From \mr{ढोंग}, a sham: \mr{ढोंग} \mr{करणे}, to put on an act. Unchanging.

\begin{cousinweb}
\textbf{\mr{ढोंगी}} (\emph{ḍhoṅgī}) is a Marathi adjective built on \mr{ढोंग}, \textquotedblleft{}a sham\textquotedblright{}, whose origin is not settled.
\end{cousinweb}

\begin{grammarlens}[title={What you've built}]
One more word for describing what you lived and felt.
\end{grammarlens}

\subsection*{Your turn}
\begin{itemize}
\raggedright
  \item \emph{Say it:} \emph{ḍhoṅgī}
  \item \emph{Say it:} \emph{ḍhoṅgī}, once more
  \item \emph{Say it:} \emph{bāher goḍ bolṇārā to māṇūs ḍhoṅgī āhe.} — That man who talks sweetly in public is a hypocrite.
  \item \emph{Recall:} say the Marathi for deceitful, then the Marathi for sly, then say \emph{ḍhoṅgī} again
\end{itemize}

\begin{tcolorbox}[breakable,colback=teal!4,colframe=teal!35!black,title={Before you move on}]
What does \mr{ढोंगी} mean? (\textquotedblleft{}Hypocritical\textquotedblright{}.) Say it once more.
\end{tcolorbox}
\section[\texorpdfstring{\mr{भोळा} (bhoḷā)}{भोळा (bhoḷā)}]{\mr{भोळा} (bhoḷā) — naive; simple-hearted}

\label{lesson:MR-C470-bhola}



\begin{quote}
Before the new one: say the Marathi for sly, then the Marathi for hypocritical.
\end{quote}

\subsection*{You'll want to know: \mr{भोळा}}
\textbf{\mr{भोळा}} — \emph{bhoḷā} — \textquotedblleft{}naive; simple-hearted\textquotedblright{}.

\begin{quote}
\mr{भोळ्या} \mr{आजोबांना} \mr{कोणीही} \mr{फसवतं}.
\end{quote}

\begin{quote}
\emph{bhoḷyā ājobānnā koṇīhī phasavtaṁ.}
\end{quote}

\begin{quote}
Anyone can fool our simple-hearted grandfather.
\end{quote}

Agrees: \mr{भोळा}, \mr{भोळी}, \mr{भोळं}. \mr{भोळाभाबडा}, guileless.

\begin{cousinweb}
\textbf{\mr{भोळा}} (\emph{bhoḷā}) is inherited: CDIAL lists Marathi \emph{bhoḷā}, \textquotedblleft{}foolish\textquotedblright{}, under the reconstructed \emph{*bhulla}, \textquotedblleft{}defective\textquotedblright{} (CDIAL 9539).
\end{cousinweb}

\begin{grammarlens}[title={What you've built}]
One more word for describing what you lived and felt.
\end{grammarlens}

\subsection*{Your turn}
\begin{itemize}
\raggedright
  \item \emph{Say it:} \emph{bhoḷā}
  \item \emph{Say it:} \emph{bhoḷā}, once more
  \item \emph{Say it:} \emph{bhoḷyā ājobānnā koṇīhī phasavtaṁ.} — Anyone can fool our simple-hearted grandfather.
  \item \emph{Recall:} say the Marathi for sly, then the Marathi for hypocritical, then say \emph{bhoḷā} again
\end{itemize}

\begin{tcolorbox}[breakable,colback=teal!4,colframe=teal!35!black,title={Before you move on}]
What does \mr{भोळा} mean? (\textquotedblleft{}Naive; simple-hearted\textquotedblright{}.) Say it once more.
\end{tcolorbox}
\section[\texorpdfstring{\mr{निरागस} (nirāgas)}{निरागस (nirāgas)}]{\mr{निरागस} (nirāgas) — innocent}

\label{lesson:MR-C470-niragas}



\begin{quote}
Before the new one: say the Marathi for hypocritical, then the Marathi for naive.
\end{quote}

\subsection*{You'll want to know: \mr{निरागस}}
\textbf{\mr{निरागस}} — \emph{nirāgas} — \textquotedblleft{}innocent\textquotedblright{}.

\begin{quote}
\mr{बाळाचा} \mr{चेहरा} \mr{किती} \mr{निरागस} \mr{आहे}!
\end{quote}

\begin{quote}
\emph{bāḷācā cehrā kitī nirāgas āhe!}
\end{quote}

\begin{quote}
How innocent the baby's face is!
\end{quote}

Unchanging; \mr{निरागसपणा}, innocence.

\begin{cousinweb}
\textbf{\mr{निरागस}} (\emph{nirāgas}) is a tatsama, taken from Sanskrit \emph{nirāgas}, \textquotedblleft{}sinless, innocent\textquotedblright{} (Monier-Williams).
\end{cousinweb}

\begin{grammarlens}[title={What you've built}]
One more word for describing what you lived and felt.
\end{grammarlens}

\subsection*{Your turn}
\begin{itemize}
\raggedright
  \item \emph{Say it:} \emph{nirāgas}
  \item \emph{Say it:} \emph{nirāgas}, once more
  \item \emph{Say it:} \emph{bāḷācā cehrā kitī nirāgas āhe!} — How innocent the baby's face is!
  \item \emph{Recall:} say the Marathi for hypocritical, then the Marathi for naive, then say \emph{nirāgas} again
\end{itemize}

\begin{tcolorbox}[breakable,colback=teal!4,colframe=teal!35!black,title={Before you move on}]
What does \mr{निरागस} mean? (\textquotedblleft{}Innocent\textquotedblright{}.) Say it once more.
\end{tcolorbox}
\section[\texorpdfstring{\mr{ठकवणे} (ṭhakavṇe)}{ठकवणे (ṭhakavṇe)}]{\mr{ठकवणे} (ṭhakavṇe) — to dupe; to swindle}

\label{lesson:MR-C470-thakavne}



\begin{quote}
Before the new one: say the Marathi for naive, then the Marathi for innocent.
\end{quote}

\subsection*{You'll want to know: \mr{ठकवणे}}
\textbf{\mr{ठकवणे}} — \emph{ṭhakavṇe} — \textquotedblleft{}to dupe; to swindle\textquotedblright{}.

\begin{quote}
\mr{एका} \mr{माणसाने} \mr{आजींना} \mr{ठकवलं} \mr{आणि} \mr{पैसे} \mr{नेले}.
\end{quote}

\begin{quote}
\emph{ekā māṇsāne ājīnnā ṭhakavlaṁ āṇi paise nele.}
\end{quote}

\begin{quote}
A man swindled Grandmother and took her money.
\end{quote}

Past, with -\mr{ने} (he …): \mr{त्याने} \mr{मला} \mr{ठकवलं}, \emph{tyāne malā ṭhakavlaṁ}.

The one cheated takes -\mr{ला}; \mr{ठक}, a swindler.

\begin{cousinweb}
\textbf{\mr{ठकवणे}} (\emph{ṭhakavṇe}) is built on the verb that CDIAL lists as Marathi \emph{ṭhakṇẽ}, \textquotedblleft{}to be deceived\textquotedblright{}, under the reconstructed \emph{*ṭhagg} (CDIAL 5489), with the causative ending \emph{-avṇe}: to make someone be deceived.
\end{cousinweb}

\begin{grammarlens}[title={What you've built}]
One more word for describing what you lived and felt.
\end{grammarlens}

\subsection*{Your turn}
\begin{itemize}
\raggedright
  \item \emph{Say it:} \emph{ṭhakavṇe}
  \item \emph{Say it:} \emph{ṭhakavṇe}, once more
  \item \emph{Say it:} \emph{ekā māṇsāne ājīnnā ṭhakavlaṁ āṇi paise nele.} — A man swindled Grandmother and took her money.
  \item \emph{Recall:} say the Marathi for naive, then the Marathi for innocent, then say \emph{ṭhakavṇe} again
\end{itemize}

\begin{tcolorbox}[breakable,colback=teal!4,colframe=teal!35!black,title={Before you move on}]
What does \mr{ठकवणे} mean? (\textquotedblleft{}To dupe; to swindle\textquotedblright{}.) Say it once more.
\end{tcolorbox}
